\documentclass[12pt,a4paper]{article}

% ─── Packages ───
\usepackage[margin=2.5cm]{geometry}
\usepackage{graphicx}
\usepackage{hyperref}
\usepackage{listings}
\usepackage{xcolor}
\usepackage{booktabs}
\usepackage{tabularx}
\usepackage{float}
\usepackage{amsmath}
\usepackage{tikz}
\usetikzlibrary{shapes.geometric, arrows.meta, positioning, fit}
\usepackage{fancyhdr}
\usepackage{titlesec}
\usepackage{enumitem}
\usepackage{caption}
\usepackage{subcaption}
\usepackage[T1]{fontenc}
\usepackage{lmodern}

% ─── Colors ───
\definecolor{codeblue}{HTML}{2196F3}
\definecolor{codegreen}{HTML}{4CAF50}
\definecolor{codered}{HTML}{F44336}
\definecolor{codegray}{HTML}{9E9E9E}
\definecolor{codebg}{HTML}{F5F5F5}
\definecolor{teal}{HTML}{00BFA5}
\definecolor{navy}{HTML}{0A1628}

% ─── Header/Footer ───
\pagestyle{fancy}
\fancyhf{}
\fancyhead[L]{\footnotesize\textsc{SCS2314 --- Middleware Architecture}}
\fancyhead[R]{\footnotesize\textsc{Presentation Script --- SwiftLogistics}}
\fancyfoot[L]{\footnotesize University of Colombo School of Computing}
\fancyfoot[R]{\footnotesize\thepage}
\renewcommand{\headrulewidth}{0.4pt}
\renewcommand{\footrulewidth}{0.2pt}

% ─── Section formatting ───
\titleformat{\section}{\Large\bfseries\color{navy}}{\thesection}{0.8em}{}[\titlerule]
\titleformat{\subsection}{\large\bfseries\color{navy!80}}{\thesubsection}{0.6em}{}
\titleformat{\subsubsection}{\normalsize\bfseries\color{navy!60}}{\thesubsubsection}{0.5em}{}

% ─── Hyperlink style ───
\hypersetup{
  colorlinks=true,
  linkcolor=codeblue,
  urlcolor=teal,
  citecolor=codegreen,
  pdftitle={SwiftTrack --- Presentation Script},
  pdfauthor={R K K Jinendra, G C K S Gajanayake, J T D Jayakodi},
  pdfsubject={SCS2314 Assignment 4}
}

% ─── Custom Environments ───
\newenvironment{scriptquote}
{\begin{quote}\itshape\color{navy!90}}
{\end{quote}}

\newcommand{\slideinfo}[1]{\noindent\textbf{\color{teal}[#1]}\medskip}

\begin{document}

% ─── Title Page ───
\begin{titlepage}
  \centering
  \vspace*{1.5cm}
  
  {\large\textsc{University of Colombo School of Computing}}\\[0.3cm]
  {\normalsize\textsc{Faculty of Science}}\\[0.2cm]
  \rule{\textwidth}{0.5pt}\\[1.5cm]
  
  {\LARGE\textsc{SCS2314 --- Middleware Architecture}}\\[0.4cm]
  {\large\textsc{Assignment 4}}\\[1.5cm]
  
  {\Huge\bfseries\color{navy} SwiftTrack}\\[0.5cm]
  {\LARGE Presentation Script for\\Screencast Walkthrough}\\[0.3cm]
  {\large\textit{SCS2314 --- Middleware Architecture Assignment 4}}\\[2cm]
  
  \rule{0.6\textwidth}{0.3pt}\\[1cm]
  
  {\large\textbf{Group Members}}\\[0.5cm]
  
  \begin{tabular}{|l|l|}
    \hline
    \textbf{Index Number} & \textbf{Name} \\
    \hline
    23000821 & R K K Jinendra \\
    \hline
    23000481 & G C K S Gajanayake \\
    \hline
    23000694 & J T D Jayakodi \\
    \hline
  \end{tabular}
  
  \vfill
  
  \begin{tabular}{rl}
    \textbf{Module:} & SCS2314 --- Middleware Architecture \\
    \textbf{Assignment:} & Assignment 4 (Middleware Application) \\
    \textbf{Time Duration:} & Under 10 Minutes \\
    \textbf{Format:} & 6-Person Narrated Screencast \\
  \end{tabular}
  
  \vspace{1cm}
  \rule{\textwidth}{0.5pt}
\end{titlepage}

\newpage

% ─── Overview ───
\section*{Script Overview}
\begin{itemize}
    \item \textbf{Total Time:} Under 10 minutes
    \item \textbf{Target Audience:} Lecturers and Peers
    \item \textbf{Language:} Simple English for clear communication
\end{itemize}

\section{Speaker 1 --- Introduction (1.5 min)}

\slideinfo{Slide: Project Title and Team Names}

\begin{scriptquote}
"Hello everyone. We are presenting \textbf{SwiftTrack}. This is our group project for the Middleware Architecture course.

Our team members are:
\begin{itemize}
    \item R K K Jinendra (23000821)
    \item G C K S Gajanayake (23000481)
    \item J T D Jayakodi (23000694)
    \item [Speaker 4 Name]
    \item [Speaker 5 Name]
    \item [Speaker 6 Name]
\end{itemize}
"
\end{scriptquote}

\slideinfo{Slide: The Problem}

\begin{scriptquote}
"Swift Logistics is a growing company. They face a big problem: they use three old systems that cannot talk to each other.

\begin{enumerate}
    \item The \textbf{Client System} uses an old SOAP format.
    \item The \textbf{Routing System} uses a modern REST format.
    \item The \textbf{Warehouse System} uses a special TCP message format.
\end{enumerate}

Our goal was to build a 'middle' system that connects all these different parts so they can work together smoothly."
\end{scriptquote}

\section{Speaker 2 --- Our Architecture (2 min)}

\slideinfo{Slide: Architecture Diagram}

\begin{scriptquote}
"To solve this, we built a microservices system. We use \textbf{RabbitMQ} as our main messenger.

This is how it works:

\begin{itemize}
    \item We have a \textbf{Client Website} where users log in.
    \item The \textbf{API Gateway} checks if the user is allowed to enter.
    \item The \textbf{Orchestrator} is the brain. It saves order details in a \textbf{MongoDB} database.
\end{itemize}

We also built three \textbf{Adapters}. Think of these as 'translators'.

\begin{itemize}
    \item One translates JSON to SOAP for the Client System.
    \item One talks to the Routing System using REST.
    \item One uses special TCP connections to talk to the Warehouse.
\end{itemize}

Finally, we have a \textbf{Notification Service}. It sends real-time updates to the website using WebSockets."
\end{scriptquote}

\section{Speaker 3 --- Why we chose this (1.5 min)}

\slideinfo{Slide: Alternative Ideas}

\begin{scriptquote}
"We looked at other ways to build this before choosing our final design.

\textbf{First Idea:} We thought about using a big central 'Bus' (ESB) like WSO2. We didn't choose this because it’s a single point of failure. If the Bus breaks, everything stops. It is also very complex to maintain.

\textbf{Second Idea:} We thought about using 'Serverless' tools like AWS Lambda. We didn't choose this because it locks us into one company like Amazon. Also, it’s not good for the Warehouse system which needs a constant connection.

\textbf{Our Choice:} We chose Microservices and RabbitMQ because it's open-source, easy to grow, and very reliable."
\end{scriptquote}

\section{Speaker 4 --- Handling Transactions (1.5 min)}

\slideinfo{Slide: Saga Flow Diagram}

\begin{scriptquote}
"A big challenge is making sure an order goes through all three systems correctly. If one part fails, we can't just leave things half-done.

We used the \textbf{Saga Pattern}. Our system follows a clear path: first CMS, then Routing, then Warehouse.

If something goes wrong at the last step, the system automatically 'undoes' the previous steps. For example, if the Warehouse is full, the system will automatically cancel the route and the inventory reservation.

This way, the data is always correct and no order is ever lost."
\end{scriptquote}

\section{Speaker 5 --- Security \& Discovery (1 min)}

\slideinfo{Slide: Security \& Service List}

\begin{scriptquote}
"Security is very important to us. We use \textbf{JWT tokens} so only logged-in users can see their data. We also have a 'Rate Limit' to stop anyone from attacking our server with too many requests.

We also added a \textbf{Service Registry}. It’s like a phonebook for our services. It checks every 10 seconds if each part of the system is 'healthy'. If a service stops working, the registry marks it as DOWN so the system knows not to use it."
\end{scriptquote}

\section{Speaker 6 --- Live Demo (2.5 min)}

\slideinfo{Action: Show the website in browser}

\begin{scriptquote}
"Now, let’s see the system in action. I am logging in as a client.

You can see the \textbf{Dashboard} now. It shows order stats like 'Pending' and 'Delivered'.

In the top corner, it says \textbf{'Connected'}. This means we are getting live updates.

Let's create a \textbf{New Order}. I will add some items and a delivery address... and click Submit.

Look! The order appears in the list. Behind the scenes, the Orchestrator talked to the CMS, the Routing system, and the Warehouse.

On the driver's side, they can see their route and mark a package as 'Delivered'. They can even take a photo or a signature as proof.

Everything works together perfectly. Thank you!"
\end{scriptquote}

\newpage
\section{Script Summary}
\begin{itemize}
    \item \textbf{Simple Words:} Replaced complex jargon like 'heterogeneous' with 'different', and 'orchestration' with 'brain' or 'coordinator' to ensure clarity.
    \item \textbf{Structure:} Divided into 6 parts (approx. 1-2 minutes per speaker) to involve all team members.
    \item \textbf{Goal:} Focuses on both the mechanical 'how' and the architectural 'why' of the middleware solution.
\end{itemize}

\end{document}
